\documentclass{ctexart}

% === 页面布局 ===
\usepackage[left=2.5cm, right=2.5cm, top=2.5cm, bottom=2.5cm]{geometry}

% === 数学与符号 ===
\usepackage{amsmath, amssymb}

% === 链接 ===
\usepackage[colorlinks,linkcolor=black,citecolor=black,urlcolor=black]{hyperref}

% === 定理类环境 ===
\newtheorem{theorem}{定理}[section]
\newtheorem{lemma}[theorem]{引理}
\newtheorem{corollary}[theorem]{推论}
\newtheorem{definition}[theorem]{定义}
\newtheorem{conjecture}[theorem]{猜想}
\newtheorem{axiom}[theorem]{公理}
\newtheorem{remark}[theorem]{注记}
\newtheorem{proof}[theorem]{证明}

\begin{document}
	
	\title{\bfseries 虚部验证机、本质维度与 $\mathbf{P} \neq \mathbf{NP}$ 的证明}
	\author{郑波尽\thanks{中南民族大学计算机学院，武汉，430074，中国。}
		\and 郑景文\thanks{武汉大学网络安全学院，武汉，430072，中国。}
		\and 王伟武\thanks{国家开发银行信息科技处，北京，100033，中国。}}
	\date{2026年7月}
	
	\maketitle
	
	\begin{abstract}
		本文在复布尔图灵机（CBTM）和实布尔图灵机（RBTM）所建立的代数语义框架基础上，完成了 $\mathbf{P} \neq \mathbf{NP}$ 的严格证明。证明分为三个逻辑上紧密关联的部分。
		
		第一部分引入虚部验证机（IVM），作为CBTM的语义展开器。CBTM/RBTM框架将非确定性刻画为代数域扩张，但其代数结构是静态的：无论发生多少次非确定性分支，所有分支共享同一对生成元，导致多次分支的代数维度坍缩为常数，对偶性也坍缩为单一的对偶对。IVM克服了这一局限。其核心创新在于动态生成元管理：借助生成元无关性定理所赋予的自由，为每次非确定性分支分配一个独立的、与之前所有生成元线性无关的素数平方根生成元，同时对偶地引入其负元。分支-激活引理确立了每次非确定性分支必然激活恰好一个生成元，这一激活是IVM作为语义展开器的定义性行为，独立于底层转移函数的设计，无法被任何算法绕过。IVM同时显式保留了每次分支的对偶性——$\sqrt{p_i}$与$-\sqrt{p_i}$的对偶关系被独立记录——修正了经典NP定义中对偶性被量词消除的错误，也从根本上消除了经典定义中隐含的虚部到实部拷贝操作。基于IVM，我们定义了本质维度 $\kappa(L)$——验证语言 $L$ 所需的最小最坏情况代数维数——并严格证明了它的良定性、实现无关性以及多项式上界。P类零维定理确立了所有 $\mathbf{P}$ 语言的本质维度为零，对应``高效可解的问题无需本质的代数扩张''。
		
		第二部分系统论证了 $\kappa(L)$ 是一个障碍无关的不变量，不受相对化、自然证明和代数化三大障碍的约束。针对相对化障碍，投影约束机制强制谕示响应编码为虚部为零的符号，将谕示限制在实部，确保 $\kappa(L)$ 在任何谕示环境下保持稳定。针对自然证明障碍，$\kappa(L)$ 的定义域严格限于 $\mathbf{NP}$ 语言，对随机布尔函数几乎处处无定义，从而结构性规避了``大尺度''条件。针对代数化障碍，非确定性分支的触发由阈值判断 $\mathfrak{Im}(\tau) > 0.5$ 决定，其核心在于区分虚部的正负号——$\sqrt{p_i}$ 与 $-\sqrt{p_i}$ 满足完全相同的代数方程，任何代数化谕示无法区分它们，因此代数化谕示对分支触发的关键信息天然``失明''。基于伽罗瓦自同构敏感性的分离证明是谕示无关证明，不受代数化障碍约束。
		
		第三部分完成 $\mathbf{P} \neq \mathbf{NP}$ 的证明。通过信息论下界论证，我们证明子集和问题的任何正确验证器都必须经历至少 $n$ 次非确定性分支：区分 $2^n$ 种可能性至少需要 $2^n$ 条不同的计算路径，而 $m$ 次二元分支至多产生 $2^m$ 条不同路径，故 $m \ge n$。结合分支-激活引理和素数平方根线性无关定理，本质维度下界为 $\Omega(n)$。验证敏感性定理从互补角度独立证明了生成元激活是验证正确性的必要条件，不可消除：构造专门实例 $(S, t=s_i)$（元素互异），唯一解迫使接受路径走特定分支，单点翻转必然改变结果。通过维度保持归约定理，$\Omega(n)$ 的维度下界传递至所有 $\mathbf{NP}$ 完全问题。反证法导出最终结论：若 $\mathbf{P} = \mathbf{NP}$，则子集和问题的本质维度同时为零和 $\Omega(n)$，矛盾。故 $\mathbf{P} \neq \mathbf{NP}$。
		
		该证明完全建立在障碍无关的不变量 $\kappa(L)$ 之上，不受相对化、自然证明和代数化三大障碍的约束。整个证明在ZFC公理体系内完成。我们进一步讨论了该证明与经典P vs.\ NP问题的关系：$\mathbf{P} \neq \mathbf{NP}$ 在集合论意义下为真，且在ZFC内可证；但该证明无法在经典图灵机模型的操作语义内部被复制，因为经典语言缺乏正面定义非确定性独立性所需的概念资源（不可公度性作为操作语义的内在属性），且经典NP定义在操作上不自足、在语义上残缺——将见证从虚部拷贝到实部（不可实现），导致不可公度性和对偶性的丢失——使得本质维度$\kappa(L)$在经典模型中根本不可定义。这一双重定位——``解决''（确认了 $\mathbf{P} \neq \mathbf{NP}$ 的真值）与``消解''（揭示了经典框架的表达局限和概念错误）并存——精确地反映了本文工作的数学内容和哲学意义。
		
		\vspace{0.5cm}
		\noindent \textbf{关键词：} 虚部验证机，本质维度，障碍无关性，分支-激活引理，验证敏感性定理，信息论下界，对偶性，操作不自足，$\mathbf{P} \neq \mathbf{NP}$
	\end{abstract}
	
	\section{引言}
	
	\subsection{P vs.\ NP问题与三大障碍}
	
	$\mathbf{P}$ versus $\mathbf{NP}$问题\cite{cook1971,levin1973}是理论计算机科学乃至整个现代数学中最具影响力的未解难题之一。$\mathbf{P} \subseteq \mathbf{NP}$是平凡的，但$\mathbf{P}$是否等于$\mathbf{NP}$——即所有可以在多项式时间内被验证的问题是否也可以在多项式时间内被求解——至今未解。若$\mathbf{P} = \mathbf{NP}$，则现代密码体系将面临严峻挑战；若$\mathbf{P} \neq \mathbf{NP}$，则将确认计算世界中存在内在的复杂性层级。
	
	经过半个多世纪的努力，研究者们逐渐识别出三个经典的证明障碍。Baker、Gill和Solovay~\cite{baker1975}的相对化障碍证明，存在两个谕示世界，一个使$\mathbf{P}=\mathbf{NP}$，另一个使$\mathbf{P}\neq\mathbf{NP}$，任何对谕示``透明''的证明技术都无法同时兼容这两个世界。Razborov和Rudich~\cite{razborov1997}的自然证明障碍指出，任何同时满足``有用性''、``大尺度''和``构造性''的组合下界证明将导致伪随机数生成器的攻破，因而被认为几乎不可能。Aaronson和Wigderson~\cite{aaronson2009}的代数化障碍将相对化推广为代数化，证明任何对谕示多项式扩展保持有效的证明技术同样无法分离$\mathbf{P}$和$\mathbf{NP}$。
	
	这些障碍共同指向一个深层问题：传统计算模型（图灵机、布尔电路）是``语义贫乏''的——它们只关心符号的语法操作，丢弃了符号背后的数学语义。要克服这些障碍，必须引入一种\textbf{基于语义的、内蕴的复杂性度量}。
	
	\subsection{从CBTM到IVM：动态维度追踪的必要性}
	
	论文\cite{models2026}已经为克服上述语义贫乏迈出了关键一步。该文基于语法不变性原理，严格证明了不可公度性不能是经典图灵机操作语义的内在属性（不可公度性元定理），并诊断了经典NP定义的操作不自足与三重概念错误——经典NP定义隐式要求了一个不可实现的虚部到实部拷贝操作；由此导致不可公度性的丢失（函数指针被扭曲为数据时抹除）、对偶性的消除（编码和量词）以及语义外包。基于这一诊断，该文构建了复布尔图灵机（CBTM）和实布尔图灵机（RBTM）。CBTM的核心洞见在于：非确定性计算对应于代数域的扩张——当读取虚部为$1$的符号时，计算必须分叉为两条路径，分别绑定于在基域$\mathbb{F}_2$上不可公度且互为对偶的生成元$\alpha$和$\beta$。这一设计将不可公度性内建为操作语义的内在属性，显式保留了对偶性，从根本上消除了拷贝操作的必要性，以语义内建取代了语义外包，系统性地修正了经典NP定义的操作不自足和三重概念错误。参数化等价定理确立了CBTM与经典模型的外延等价（$\mathbf{P}_{cb}=\mathbf{P}$，$\mathbf{NP}_{cb}=\mathbf{NP}$）。RBTM进一步引入了双带视角和生成元无关性定理，揭示了非确定性的本质在于``引入与基域不可公度的新元素''这一事实本身，并揭示了伽罗瓦自同构的操作性本质——路径切换——据此提出了与$\mathbf{P}\neq\mathbf{NP}$等价的敏感性猜想。
	
	然而，CBTM/RBTM框架存在一个根本性的局限：其代数结构是\textbf{静态}的。无论计算过程中发生多少次非确定性分支，所有符号的虚部非零部分都只能关联于同一对生成元$\{\alpha, \beta\}$（或$\{\sqrt{2}, 1+\sqrt{2}\}$），整个计算所张成的代数空间维数最多为$2$，对偶性也坍缩为单一的对偶对。因此，CBTM虽然完美地刻画了单次分支的独立性（不可公度性）和对偶性，但无法区分``发生一次分支''与``发生了$n$次分支''——多次分支的代数维度会坍缩回一个常数。对于$\mathbf{NP}$完全问题，其验证过程本质上需要大量独立的非确定性选择，而CBTM的固定维度无法反映这种需求，因而不能直接用于分离$\mathbf{P}$与$\mathbf{NP}$完全问题。
	
	本文的核心贡献正是克服这一局限。我们引入\textbf{虚部验证机（Imaginary-part Verification Machine, IVM）}，作为CBTM的语义展开器。IVM的核心创新在于\textbf{动态生成元管理}：借助生成元无关性定理所赋予的自由，为每次非确定性分支分配一个独立的、与之前所有生成元线性无关的素数平方根生成元$\sqrt{p_i}$，同时对偶地引入$-\sqrt{p_i}$。由素数平方根线性无关定理，这些生成元彼此线性无关，从而使得代数维度能够随分支次数线性累积，对偶性也得以按分支步骤独立追踪，彻底解决了静态框架的维度坍缩和对偶性坍缩问题。值得强调的是，IVM的语义展开直接保留了虚部中的不可公度性和对偶性，无需任何从虚部到实部的拷贝——这正是对经典NP定义中操作不自足的根本修正。
	
	基于IVM，我们定义了\textbf{本质维度$\kappa(L)$}——验证语言$L$所需的最小最坏情况代数维数。我们系统论证了$\kappa(L)$是一个障碍无关的复杂性不变量，能够区分$\mathbf{P}$与$\mathbf{NP}$而不受三大障碍干扰。利用$\kappa(L)$，我们最终完成了$\mathbf{P} \neq \mathbf{NP}$的严格证明：P类零维定理确立了所有$\mathbf{P}$语言的本质维度为零；通过信息论下界论证与分支-激活引理，证明了子集和问题的本质维度为$\Omega(n)$；反证法导出最终结论。
	
	\subsection{本文的贡献与结构}
	
	本文的主要贡献包括：
	
	\begin{enumerate}
		\item \textbf{虚部验证机（IVM）与本质维度$\kappa(L)$}（第3节）：构建了IVM作为CBTM的语义展开器，通过动态生成元管理机制解决了静态框架的维度坍缩和对偶性坍缩问题。定义了本质维度$\kappa(L)$，并严格证明了它的良定性、实现无关性、多项式上界以及P类零维定理（所有$\mathbf{P}$语言的$\kappa=0$）。分支-激活引理确立了维度消耗是IVM语义记账的强制性结果，独立于底层转移函数的设计，无法被任何算法绕过。IVM同时显式保留了每次分支的对偶性，修正了经典NP定义中对偶性被量词消除的错误。
		\item \textbf{障碍无关性论证}（第4节）：系统证明了$\kappa(L)$不受相对化、自然证明和代数化三大障碍的约束。投影约束机制确保$\kappa(L)$在任何谕示下保持稳定；定义域限制使其对随机布尔函数几乎处处无定义，规避了自然证明的``大尺度''条件；阈值判断的完全离散性使分支触发依赖于符号的正负号——即对偶生成元的选择——代数化谕示对此天然``失明''，基于伽罗瓦敏感性的分离证明因此是谕示无关证明。
		\item \textbf{NP完全问题的维度下界}（第5节）：通过信息论下界论证，证明了子集和问题的任何正确验证器都必须经历至少$n$次非确定性分支，结合分支-激活引理和素数平方根线性无关定理，得出$\kappa(\mathrm{SubsetSum}) = \Omega(n)$。验证敏感性定理独立证明了生成元激活是验证正确性的必要条件，不可消除。维度保持归约定理将$\Omega(n)$的下界传递至所有$\mathbf{NP}$完全问题。
		\item \textbf{$\mathbf{P} \neq \mathbf{NP}$的证明与经典框架的定位}（第6-7节）：通过反证法完成$\mathbf{P} \neq \mathbf{NP}$的证明，并讨论了该证明与经典P vs.\ NP问题的关系：$\mathbf{P} \neq \mathbf{NP}$在集合论意义下为真，且在ZFC内可证；但该证明无法在经典图灵机模型的操作语义内部被复制，因为经典语言缺乏正面定义非确定性独立性所需的概念资源，且经典NP定义在操作上不自足、在语义上残缺，使得本质维度$\kappa(L)$在经典模型中不可定义。
	\end{enumerate}
	
	本文的结构如下。第2节回顾必要的预备知识，包括CBTM/RBTM框架的核心定义和定理。第3节构建IVM并定义本质维度$\kappa(L)$。第4节论证$\kappa(L)$的障碍无关性。第5节证明NP完全问题的维度下界。第6节完成$\mathbf{P} \neq \mathbf{NP}$的证明。第7节讨论结果对经典框架的意义。第8节总结全文。
	
	\section{预备知识}
	
	\subsection{CBTM/RBTM框架回顾}
	
	本文的整个技术框架建立在论文\cite{models2026}所构建的复布尔图灵机（CBTM）和实布尔图灵机（RBTM）的基础上。为保持本文的自足性，本节简要回顾CBTM和RBTM的核心定义和结论，详细内容请参见原文。
	
	\subsubsection*{复布尔图灵机（CBTM）}
	
	CBTM的纸带符号取自四元伽罗瓦域$\mathbb{F}_4 = \{0, 1, \alpha, \beta = \alpha+1\}$，其中$\mathbb{F}_2 = \{0,1\}$为基域。投影算子$\operatorname{Re}, \operatorname{Im}: \mathbb{F}_4 \to \{0,1\}$将每个符号分解为实部和虚部：$\operatorname{Re}(0)=\operatorname{Re}(\alpha)=0$，$\operatorname{Re}(1)=\operatorname{Re}(\beta)=1$；$\operatorname{Im}(0)=\operatorname{Im}(1)=0$，$\operatorname{Im}(\alpha)=\operatorname{Im}(\beta)=1$。代数语义上，每个符号$s$可写为$s = \operatorname{Re}(s) + \operatorname{Im}(s) \cdot \alpha$。$\operatorname{Im}(s)=1$意味着该符号涉及域扩张中的不可公度元$\alpha$。
	
	CBTM由两条核心公理支配。\textbf{分支触发公理}：当读取符号的虚部$\operatorname{Im}(\tau)=1$时，转移关系恰好提供两个后继选择（非确定性模式），分别绑定于对偶生成元$\alpha$和$\beta$；当虚部为$0$时，恰好提供一个后继选择（确定性模式）。\textbf{投影约束公理}：写入符号的虚部完全由当前状态和读取符号的虚部通过布尔函数决定；特别地，若读取符号虚部为零，则写入符号虚部也必须为零。CBTM的纸带可以可视化为两条同步的布尔纸带——实带（确定性数据）和虚带（非确定性控制）。
	
	CBTM与经典计算模型的关系由\textbf{参数化等价定理}确立：限制在所有符号虚部为$0$的CBTM（记作$\mathrm{CBTM}|_0$）与经典确定性图灵机多项式等价；完整的CBTM与经典非确定性图灵机多项式等价。因此$\mathbf{P}_{cb} = \mathbf{P}$，$\mathbf{NP}_{cb} = \mathbf{NP}$。CBTM并未引入超计算能力，它只是为已有的非确定性计算提供了一种全新的、带有代数语义的解释，并系统性地修正了经典NP定义的操作不自足和三重概念错误——函数指针常驻虚部（保留了不可公度性）、对偶性显式保留、拷贝操作被消除、语义内建而非外包。参数化猜想将$\mathbf{P}$ vs.\ $\mathbf{NP}$等价转化为虚部$b$是否可消除的代数问题：域扩张$\mathbb{F}_4/\mathbb{F}_2$是否为非确定性计算提供了本质的代数资源？
	
	\subsubsection*{实布尔图灵机（RBTM）与生成元无关性}
	
	RBTM将抽象生成元$\alpha$具体化为代数数$\sqrt{2}$，提出双带视角：每条纸带被分解为实带（存储有理系数$a$，对应确定性数据）和虚带（存储无理系数$b$，对应非确定性控制）。RBTM与CBTM在计算能力上等价。
	
	论文\cite{models2026}最重要的结果是\textbf{不可公度性元定理}和\textbf{生成元无关性定理}。不可公度性元定理严格证明了经典图灵机无法内在地表达不可公度性，这一定理构成了整个维度退化理论的元理论基石。生成元无关性定理证明了计算能力不依赖于生成元的具体选择——无论使用$\sqrt{2}$、$\sqrt{3}$还是虚数单位$i$作为生成元，所构造的自动机都是同构的，识别相同的语言类。这一定理赋予了我们在IVM中自由选择最方便生成元体系的权力——我们可以使用素数平方根序列$\sqrt{p_1}, \sqrt{p_2}, \dots$来标记不同的非确定性分支步骤，而不改变计算能力。
	
	\subsection{生成元无关性与双带视角}
	
	生成元无关性定理的深层意义在于：它揭示了非确定性的本质在于``引入与基域不可公度的新元素''这一事实本身，而非生成元的个体代数身份。这使得我们在IVM中可以自由选择最方便于度量的生成元体系，而不扭曲问题的本质。
	
	双带视角将实部（数据）与虚部（控制）分离，使``维度''概念可视化。双带视角的核心贡献之一是消解了有限域$\mathbb{F}_4$与无限域$\mathbb{Q}(\sqrt{2})$之间的域非同构问题：尽管两者作为代数结构并不同构（特征不同），但它们在布尔投影层所诱导的自动机却是同构的。这意味着计算语义的同一性——包括不可公度性和对偶性——可以超越底层域结构的差异。
	
	\subsection{素数平方根线性无关定理}
	
	IVM能够分离$\mathbf{P}$与$\mathbf{NP}$完全问题的数学基石是以下经典结果。
	
	\begin{lemma}[素数平方根线性无关定理]\label{lem:commensurability}
		设$p_1, \dots, p_m$为互异素数。则集合$\{\sqrt{p_1}, \dots, \sqrt{p_m}\}$在$\mathbb{Q}$上线性无关。
	\end{lemma}
	
	该定理的证明可参见代数数论的标准教材\cite{lang2002}。它保证了当IVM为每次非确定性分支分配一个独立的素数平方根生成元时，这些生成元彼此线性无关，从而使得代数维度能够随分支次数线性累积，忠实地反映非确定性分支的次数。没有这一定理，IVM将无法将多次分支的维度``计数''为线性增长——这正是CBTM静态框架所面临的维度坍缩问题的反面。
	
	\subsection{敏感性猜想与路径切换判据}
	
	论文\cite{models2026}揭示了伽罗瓦自同构的操作性本质——路径切换——并据此提出了敏感性猜想和路径切换判据。这些结果与本文的论证密切相关。
	
	伽罗瓦自同构$\sigma: \alpha \leftrightarrow \beta$在计算树上的作用是将每条路径的所有非确定性选择全部翻转，即切换到其对偶路径。在IVM的动态扩张下，伽罗瓦群从单一的$\mathbb{Z}/2\mathbb{Z}$扩展为$(\mathbb{Z}/2\mathbb{Z})^m$，允许独立翻转单个非确定性选择点（单点敏感性）。
	
	\textbf{单点敏感性}：机器$M$在输入$x$上对第$i$个非确定性选择点敏感，如果存在一条接受路径$\pi$，使得单独翻转$\pi$上第$i$个选择后，所得路径不再是接受路径。若$M$在所有输入上对所有选择点均不敏感，则称$M$\textbf{完全不敏感}。
	
	\textbf{路径切换判据}：若语言$L$可由一台完全不敏感的CBTM/RBTM判定，则$L \in \mathsf{P}$。若$\mathsf{P} \neq \mathsf{NP}$，则对任何$L \in \mathsf{NP} \setminus \mathsf{P}$，任何接受$L$的机器在某些输入上必然存在某个非确定性选择点，使得机器在该点上单点敏感。特别地，NP完全语言在任何接受它的机器上都具有不可消除的单点敏感性。
	
	\textbf{敏感性猜想}：存在一个NP完全语言，使得任何判定它的多项式时间CBTM（或RBTM）都必然在某个输入上对至少一个非确定性选择点单点敏感。该猜想等价于$\mathbf{P}\neq\mathbf{NP}$。
	
	敏感性猜想将P vs NP问题转化为一个定性结构判据：非确定性计算的力量是否必然导致存在某个本质的非确定性选择点，单独翻转其选择就会改变接受状态？这一翻转操作在代数上对应于将对偶生成元的选择从 $+\sqrt{p_i}$ 切换为 $-\sqrt{p_i}$（或反之），即对偶性在操作层面的直接体现。在本文的IVM框架中，每一个本质的非确定性选择点恰好对应一个被激活的生成元，因此本质维度$\kappa(L)$恰好度量了本质选择点的数量。敏感性猜想的成立将直接导出$\kappa(L) = \Omega(n)$的下界，从而完成$\mathsf{P} \neq \mathsf{NP}$的证明。
	
	
	
	\section{虚部验证机与本质维度}
	
	本节引入虚部验证机（IVM）作为CBTM的语义展开器，解决静态框架的维度坍缩和对偶性坍缩问题，并在此基础上定义本质维度$\kappa(L)$。与CBTM的静态代数结构不同，IVM通过动态生成元管理机制，使得代数维度能够随非确定性分支的次数线性累积，对偶性也得以按分支步骤独立追踪，从而为度量非确定性资源消耗提供了精确的代数工具。
	
	\subsection{IVM的概念定位：面向NPC度量的语义展开器}
	
	IVM的目的并非执行计算，而是\textbf{记录与分析}CBTM的计算过程，尤其是那些涉及大量非确定性分支的验证过程。其输入为一台CBTM $M$和一个输入$x$，输出则是$M$在$x$上的完整计算轨迹，并以代数语义的形式呈现。IVM本质上是对CBTM计算的一个子集——即分支次数随输入规模增长、需要独立记录每次分支选择的计算——进行语义展开。对于纯确定性计算或分支次数为常数的计算，CBTM自身的结构已足够；IVM的真正价值体现在$\mathbf{NP}$完全问题的验证过程中，正是在这些场景下，区分不同分支步骤所引入的维度变得至关重要。
	
	我们可以将IVM理解为CBTM的``度量强化版''：CBTM首次将非确定性从语法约定提升为代数事实，内建了不可公度性，显式保留了对偶性，消除了拷贝操作，但其代数结构是静态的——无论发生多少次非确定性分支，所有分支共享同一对生成元，导致维度坍缩为常数，对偶性也坍缩为单一的对偶对。IVM通过动态生成元管理，使得代数维度能够随分支次数线性累积，对偶性也得以按分支步骤独立追踪（每次分支分配独立的对偶对$\sqrt{p_i}$与$-\sqrt{p_i}$），从而能够定量区分``发生了一次分支''与``发生了$n$次分支''在代数资源消耗上的本质差异。
	
	\subsection{IVM的形式化定义与格局}
	
	\begin{definition}[IVM格局]\label{def:ivm-config}
		在第$t$步的IVM格局定义为一个五元组
		\[
		C_t = (q_t, T_t^{\mathrm{Re}}, T_t^{\mathrm{Im}}, n_t, h_t),
		\]
		其中：
		\begin{itemize}
			\item $q_t\in Q$：控制状态，与底层CBTM的当前状态相同；
			\item $T_t^{\mathrm{Re}}:\mathbb{Z}\to\{0,1\}$：实带，存储每个位置的实系数；
			\item $T_t^{\mathrm{Im}}:\mathbb{Z}\to\{0,1\}$：虚系数带，存储每个位置的虚系数；
			\item $n_t\in\mathbb{N}$：生成元计数器，记录到第$t$步为止已引入并激活的代数生成元的个数；
			\item $h_t\in\mathbb{Z}$：带头位置。
		\end{itemize}
	\end{definition}
	
	实带与虚系数带共同决定了位置$i$处的完整代数符号：$\tau_i = T_t^{\mathrm{Re}}[i] + T_t^{\mathrm{Im}}[i]\cdot\sqrt{p_{n_i}}$。这里$\sqrt{p_{n_i}}$是第$n_i$个素数平方根生成元，其对偶为$-\sqrt{p_{n_i}}$。生成元计数器$n_t$是IVM的核心创新，它动态地追踪代数维度的演化——每当发生非确定性分支，$n_t$递增，并分配一个新的、与之前所有生成元线性无关的素数平方根，同时对偶地引入其负元。这使得IVM能够精确记录每次分支所消耗的代数资源，并独立追踪每次分支的对偶结构。
	
	\subsection{动态生成元管理与分支-激活引理}
	
	\begin{definition}[生成元更新规则]\label{def:generator-update}
		生成元计数器按以下规则更新：
		\[
		n_{t+1} = \begin{cases}
			n_t + 1 & \text{若 } |T_t^{\mathrm{Im}}[h_t]| > \epsilon \ (\text{触发分支}), \\
			n_t & \text{否则（确定性转移）}.
		\end{cases}
		\]
		当$n_t$增加时，新引入的生成元被分配为第$n_t$个素数平方根$\sqrt{p_{n_t}}$，其中$p_1=2,p_2=3,p_3=5,\dots$为素数序列，同时对偶地引入$-\sqrt{p_{n_t}}$。分支阈值$\epsilon$通常取$0.5$，以便将虚系数$\{0,1\}$明确分离。
	\end{definition}
	
	在IVM的语义展开中，生成元的引入与激活是同一事件的两个不可分割的侧面。当分支触发时，IVM自动执行以下操作：
	\begin{enumerate}
		\item 计数器$n_t$增加$1$，分配一个新的素数平方根地址$\sqrt{p_{n_t}}$，同时对偶地引入$-\sqrt{p_{n_t}}$；
		\item 在激活路径上记录值$1$，在不激活路径上记录值$0$。这里激活与不激活分别对应两个对偶生成元的选择——激活对应$+\sqrt{p_{n_t}}$，不激活对应$-\sqrt{p_{n_t}}$。
	\end{enumerate}
	这一标记过程是IVM作为语义展开器的定义性行为，独立于底层CBTM转移函数的设计。值$1$就是激活的标志，值$0$就是不激活的标志。因此，\textbf{在IVM的语义层面，生成元的引入即意味着其被激活}——每次分支必然激活恰好一个生成元，即激活路径上记录的值$1$所对应的素数平方根生成元。与此同时，对偶生成元$-\sqrt{p_{n_t}}$在不激活路径上被隐式保留。我们将这一核心事实总结为以下引理。
	
	\begin{lemma}[分支-激活引理]\label{lem:branch-activate}
		设$M$为任意IVM。若在某个计算步骤中发生非确定性分支，则该分支必然激活恰好一个生成元。具体而言，在分支触发的瞬间，IVM自动分配一个新的素数平方根地址$\sqrt{p_i}$（同时对偶地引入$-\sqrt{p_i}$），并在激活路径上记录值$1$，对应的素数平方根生成元被激活。这一激活是IVM作为语义展开器的定义性行为，独立于底层转移函数的设计，无法被任何算法绕过。
	\end{lemma}
	
	\begin{remark}[生成元激活的不可绕过性与对偶性的显式保留]
		分支-激活引理揭示了一个关键事实：维度消耗并非验证器的主动选择，而是IVM语义记账的强制性结果。类比于经典图灵机的步数统计无法被转移函数``关闭''——无论转移函数如何设计，每执行一步，步数计数器就增加一——IVM的生成元激活记录同样无法被任何巧妙设计的转移函数所绕过。无论底层CBTM在纸带上写入何种符号（虚部为零或非零），IVM在语义层面已经为激活路径记录了值$1$，该生成元已被激活，其对偶生成元$-\sqrt{p_i}$在不激活路径上被隐式保留。本质维度$\kappa(L)$统计的正是这些语义层面的激活记录。因此，不存在通过巧妙设计转移函数来规避维度消耗的可能性。
		
		值得注意的是，IVM同时修正了经典NP定义中对偶性被存在量词消除的错误。在IVM中，每次非确定性分支的对偶性——$\sqrt{p_i}$与$-\sqrt{p_i}$之间的伽罗瓦对称性——被显式保留在激活/不激活的布尔标记中。对偶性不再是经典定义中被抹去的隐性结构，而是IVM操作语义中可独立追踪的显式属性。与此同时，IVM的语义展开完全不涉及从虚部到实部的拷贝——分支选择直接记录在虚部地址空间中，这正是对经典NP定义操作不自足的根本修正。
	\end{remark}
	
	\textbf{两个定理的协同：从不可公度性到线性无关性。} IVM的生成元管理策略依赖于两个独立的数学定理，它们共同为分离$\mathbf{P}$与$\mathbf{NP}$完全问题提供了基础。\textbf{生成元无关性定理}\cite{models2026}保证IVM可以自由选择任意代数数作为生成元，而不改变识别语言类。这赋予了我们在不扭曲问题本质的前提下，选用最方便于度量的生成元体系的权力。\textbf{素数平方根线性无关定理}（引理\ref{lem:commensurability}）则保证当我们选用素数平方根序列时，每次分支引入并激活的新生成元与之前所有生成元线性无关。
	
	在CBTM中，不可公度性只保证了单次分支的路径独立性和对偶性，但所有分支共享同一对生成元，导致多次分支的维度坍缩为常数，对偶性也坍缩为单一的对偶对。IVM通过上述两个定理的协同，将CBTM中的``不可公度性''（两两关系）升级为``线性无关性''（整体关系），并将单一对偶对扩展为独立的对偶对序列：生成元无关性定理允许我们动态分配新生成元，而线性无关定理确保这些生成元彼此独立，从而维度得以线性累积，对偶性也得以按分支步骤独立追踪。这一跨越使得本质维度$\kappa(L)$能够精确地区分$\mathbf{P}$类语言（维度为零）与$\mathbf{NP}$完全语言（维度随输入规模线性增长）。
	
	\begin{remark}[生成元选择对计算与度量的不同影响]
		IVM的生成元分配策略独立于底层CBTM的转移函数，因此无论选择何种生成元序列，底层CBTM的计算行为（计算树结构、接受/拒绝结果）完全不受影响。生成元选择影响的仅仅是IVM对非确定性资源的度量精度。若选择线性无关的生成元序列（如素数平方根），IVM能够精确度量多次独立非确定性分支的维度累积，适用于$\mathbf{NP}$完全问题的复杂度分析；若选择线性相关的生成元序列（如重复使用同一生成元$\sqrt{2}$），IVM退化为CBTM的忠实记录器，维度坍缩为常数，对偶性坍缩为单一对偶对，适用于非$\mathbf{NP}$完全问题的记录。这种灵活性正是IVM作为统一框架的精妙之处：它既是对CBTM的忠实继承（在兼容模式下），也是对CBTM的超越（在度量模式下）。
	\end{remark}
	
	\subsection{从CBTM到IVM的语义展开}
	
	IVM作为CBTM的语义展开器，其忠实性体现在两者之间存在严格的形式化对应关系。设底层CBTM的计算过程包含$m$次非确定性分支。根据CBTM的分支触发公理和投影约束公理，每次分支的两条路径分别绑定到不可公度且互为对偶的生成元$\alpha$和$\beta$。第$i$次分支的选择记为$\pi_i \in \{\alpha, \beta\}$。整条路径的标记是长度为$m$的序列$\Pi = (\pi_1, \pi_2, \dots, \pi_m) \in \{\alpha, \beta\}^m$。
	
	在IVM中，每次分支触发时分配一个全新的素数平方根地址$\sqrt{p_i}$（同时对偶地引入$-\sqrt{p_i}$），并在激活路径上记录值$1$（激活该生成元），在不激活路径上记录值$0$。整条路径的标记是长度为$m$的布尔向量$\mathbf{a} = (a_1, a_2, \dots, a_m) \in \{0, 1\}^m$，其中$a_i = 1$表示该路径在第$i$次分支处走了激活路径（选择了$+\sqrt{p_i}$），$a_i = 0$表示走了不激活路径（选择了$-\sqrt{p_i}$）。
	
	\begin{definition}[语义展开映射]\label{def:semantic-unfolding}
		CBTM到IVM的语义展开由映射$\Phi: \{\alpha, \beta\}^m \to \{0,1\}^m$给出：
		\[
		\Phi(\Pi) = \mathbf{a}, \quad \text{其中 } a_i = 
		\begin{cases}
			1 & \text{若 } \pi_i = \alpha,\\
			0 & \text{若 } \pi_i = \beta.
		\end{cases}
		\]
	\end{definition}
	
	\begin{theorem}[语义展开是一一对应]\label{thm:phi-bijection}
		映射$\Phi$是双射。
	\end{theorem}
	\begin{proof}
		定义逆映射$\Phi^{-1}: \{0,1\}^m \to \{\alpha, \beta\}^m$为$\Phi^{-1}(\mathbf{a})_i = \alpha$若$a_i=1$，否则$\Phi^{-1}(\mathbf{a})_i = \beta$。显然$\Phi \circ \Phi^{-1} = \text{id}$且$\Phi^{-1} \circ \Phi = \text{id}$。
	\end{proof}
	
	在CBTM中，伽罗瓦自同构$\sigma: \alpha \leftrightarrow \beta$作用于路径标记$\Pi$上逐分量交换地址。在IVM中，对应的操作为路径切换$\rho: \{0,1\}^m \to \{0,1\}^m$，定义为$\rho(\mathbf{a}) = \mathbf{a} \oplus \mathbf{1}_m$。
	
	\begin{theorem}[语义展开保持伽罗瓦作用]\label{thm:phi-galois}
		$\Phi \circ \sigma = \rho \circ \Phi$。
	\end{theorem}
	\begin{proof}
		任取$\Pi \in \{\alpha, \beta\}^m$。设$a_i = 1$若$\pi_i = \alpha$，$a_i = 0$若$\pi_i = \beta$。则$\Phi(\sigma(\Pi))_i = 1-a_i = \rho(\Phi(\Pi))_i$。
	\end{proof}
	
	\textbf{维度累积与对偶性追踪的保证。} 在CBTM中，所有路径标记的分量取自同一对地址$\{\alpha, \beta\}$，因此虚部空间维数最多为$2$，对偶性也限于这一对生成元。在IVM中，地址空间$\mathcal{A}_{\text{IVM}} = \{\sqrt{p_1}, \dots, \sqrt{p_m}\}$的大小随$m$线性增长，且由素数平方根线性无关定理，这些地址彼此线性无关。每次分支激活恰好一个生成元，同时对偶生成元$-\sqrt{p_i}$在不激活路径上被隐式保留。因此$m$次分支激活的生成元张成空间的维数为$m$，对偶性也得以按分支步骤独立追踪。语义展开映射$\Phi$将CBTM中固定地址空间内的标记，转换为IVM中动态扩展地址空间上的布尔值，从而实现了维度从常数到线性的解放，以及对偶性从单一对偶对到独立对偶对序列的解放。这一解放是分离$\mathbf{P}$与$\mathbf{NP}$完全问题的关键：只有在线性维度下，我们才能观察到$\mathbf{P}$（零维）与$\mathbf{NP}$完全（无界维）之间的本质差异。
	
	\subsection{代数封闭性与计算可行性}
	
	IVM的动态域扩张必须在代数上保持封闭，且在计算上保持可行。本节给出这两方面的严格保证。
	
	\begin{theorem}[基本代数运算的封闭性]\label{thm:basic-closure}
		设$K = \mathbb{Q}(\sqrt{p_1},\dots,\sqrt{p_k})$为IVM所产生的代数数域，其中$p_1,\dots,p_k$为互异素数。则IVM的加法、乘法、数乘和线性组合运算均在$K$内封闭。
	\end{theorem}
	\begin{proof}
		封闭性直接源于$K$作为$\mathbb{Q}$的扩域（从而是一个域）的定义。域对加法、乘法封闭；有理数乘法保持在$K$中因为$\mathbb{Q}\subset K$；线性组合则是有限次加法和数乘的复合。
	\end{proof}
	
	\begin{theorem}[动态扩张下的封闭性]\label{thm:dynamic-closure}
		在IVM执行过程中，若在步骤$t$动态引入并激活新生成元$\sqrt{p_{k+1}}$，将域扩展为$K_t = \mathbb{Q}(\sqrt{p_1},\dots,\sqrt{p_k},\sqrt{p_{k+1}})$，则：
		\begin{enumerate}
			\item $[K_t:\mathbb{Q}] = 2^{k+1} < \infty$（保持有限维）；
			\item 原有运算在$K_t$中仍封闭，因为$K\subset K_t$；
			\item 新生成元的引入遵循逐步赋值策略，保证线性无关；
			\item 时间开销增加$O(\log^2 n)$，保持在多项式界内。
		\end{enumerate}
	\end{theorem}
	
	\begin{theorem}[交叉项的规避]\label{thm:cross-term}
		IVM在执行过程中不会生成新的交叉项$\sqrt{p_i p_j}$ ($i\neq j$)。所有代数运算均保持纯项结构。
	\end{theorem}
	\begin{proof}
		IVM记录底层CBTM的计算轨迹，而CBTM的每一步写入符号遵循投影约束：写入符号$\gamma_w = f_{\mathfrak{Re}}(q,\mathfrak{Re}(\tau)) + \alpha \cdot f_{\mathfrak{Im}}(q,\mathfrak{Im}(\tau))$，其中$\alpha$是当前步骤所关联的单一生成元（由IVM分配）。由于$f_{\mathfrak{Re}}, f_{\mathfrak{Im}}$取值于$\{0,1\}$，$\gamma_w$只能是$0,1,\alpha,1+\alpha$之一，均为纯项。初始纸带上的符号（输入编码）也被设计为纯项。代数域的扩张仅通过引入并激活新生成元发生，而转移函数从未产生不同生成元的乘积。故交叉项不会出现。
	\end{proof}
	
	这一规避机制确保了计算过程中涉及的代数数始终保持在可有效处理的范围内——每个符号只涉及单个素数平方根，不会出现组合爆炸。它为本质维度$\kappa(L)$的度量提供了清晰的代数基础：维度计数仅取决于被激活的生成元数量，而非生成元之间的乘积组合。
	
	\subsection{生成元提取算子与单次验证的维度}
	
	\begin{definition}[生成元提取算子]\label{def:generator-extraction}
		设$K$为IVM执行过程中最终生成的代数域。定义$\operatorname{\mathfrak{G}}:\Gamma\to K$为：
		\[
		\operatorname{\mathfrak{G}}(\tau)=\begin{cases}
			\sqrt{p_i} & \text{若 }\tau = a + b\sqrt{p_i}\in K\text{ 且 }b\neq0,\\
			0 & \text{否则}.
		\end{cases}
		\]
	\end{definition}
	
	生成元提取算子是定义$\kappa(L)$的基础。它从每个符号中提取其所涉及的生成元信息——若符号的虚部非零，则返回对应的素数平方根；否则返回$0$。在静态RBTM框架中，所有分支共享同一个生成元$\sqrt{2}$，提取算子只能返回$\sqrt{2}$或$0$，无法区分不同分支步骤。这正是静态框架的局限性，凸显了引入动态IVM的必要性。在IVM中，每个分支步骤被分配独立的素数平方根生成元，提取算子才能精确记录维度的增长。
	
	\begin{definition}[单次验证的维度]\label{def:single-dim}
		对于IVM $M$和输入$x$，定义
		\[
		\kappa_M(x) = \dim_{\mathbb{Q}}\left\langle \left\{ \operatorname{\mathfrak{G}}(\tau_t) \mid t\in\mathrm{AllPaths}(M,x)  \right\} \right\rangle_{\mathbb{Q}},
		\]
		其中$\tau_t$表示某条路径上第$t$步的符号，$\mathrm{AllPaths}(M,x)$为$M$在$x$上所有计算路径（作为配置序列）的集合。
	\end{definition}
	
	根据分支-激活引理，每次分支必然激活恰好一个生成元（在激活路径上记录值$1$）。因此，$\kappa_M(x)$正是所有被激活的生成元张成空间的维数。由于生成元的激活是IVM语义展开的强制性记账，$\kappa_M(x)$的值独立于底层CBTM转移函数的设计。这一性质使得$\kappa_M(x)$成为一个可靠的计算资源度量——它不会被任何``取巧''的转移函数设计所压低。
	
	\subsection{本质维度$\kappa(L)$的定义}
	
	为消除特定IVM实现的偏差，并捕捉语言$L$本身在最坏情况下所需的最小代数资源，我们采用``极小—极大''定义：
	
	\begin{definition}[本质维度]\label{def:kappa}
		对语言$L\in\mathbf{NP}$，定义
		\[
		\boxed{\ \kappa(L)=\inf_{M\in\mathcal{V}_L}\sup_{x\in\Sigma^*}\kappa_M(x)\ }.
		\]
		这里$\mathcal{V}_L$是所有在多项式时间内验证$L$的IVM的集合。
	\end{definition}
	
	该定义的直观含义是：$\sup_x$捕获最坏情况下的维度需求；$\inf_M$在所有可能的验证机中选取最优者，反映了语言$L$\textbf{必须消耗的、不可压缩的}维度下界。任何特定的验证机可能通过``铺张浪费''的方式引入不必要的生成元，从而得到更高的维度；但$\kappa(L)$取下确界，排除了这种人为膨胀，抓住了语言$L$的内在代数资源需求。
	
	\subsection{良定性、实现无关性与多项式上界}
	
	\begin{theorem}[良定性]\label{thm:well-defined}
		对任意$L\in\mathbf{NP}$，$\kappa(L)$存在（可以为$\infty$）。
	\end{theorem}
	\begin{proof}
		由IVM与NTM的等价性（见论文\cite{models2026}的参数化等价定理），$\mathbf{NP}$语言有多项式时间NTM，因而也有多项式时间IVM验证器，故$\mathcal{V}_L\neq\emptyset$。对每个$M\in\mathcal{V}_L$，由于$M$的运行时间多项式有界且每步至多引入并激活一个生成元，$\kappa_M(x)$有上界，从而$\sup_x\kappa_M(x)$有限。取其下确界即得$\kappa(L)$。
	\end{proof}
	
	\begin{theorem}[实现无关性]\label{thm:implementation-indep}
		对任意$M_1,M_2\in\mathcal{V}_L$，存在多项式$p$使得对一切$x\in\Sigma^*$，$|\kappa_{M_1}(x)-\kappa_{M_2}(x)|\le p(|x|)$。因而$\kappa(L)$不依赖于具体的IVM实现。
	\end{theorem}
	\begin{proof}[证明思路]
		通过NTM相互模拟：$M_1$与$M_2$可转化为NTM $N_1,N_2$，而两者可互相模拟。模拟过程中每一步引入的额外生成元数量受多项式控制，故维度差异被多项式界限所约束。详细证明见附录A。
	\end{proof}
	
	\begin{theorem}[多项式时间上界]\label{thm:poly-bound}
		存在多项式$p$，使得对任何$M\in\mathcal{V}_L$及任何输入$x$，$\kappa_M(x)\le p(|x|)$。
	\end{theorem}
	\begin{proof}
		$M$运行时间$T(|x|)$为多项式，每步至多引入并激活一个生成元，故$\kappa_M(x)\le T(|x|)$。
	\end{proof}
	
	\subsection{P类零维定理}
	
	\begin{theorem}[P类的零维实现]\label{thm:p-zero}
		若$L\in\mathbf{P}$，则$\kappa(L)=0$。
	\end{theorem}
	\begin{proof}
		$L\in\mathbf{P}$意味着存在确定性图灵机$M_{\det}$在多项式时间内判定$L$。构造IVM $M$如下：状态空间为$Q_{\det}\times\{0\}$（添加平凡的虚分量）；符号集为$\Sigma_{\det}\times\{0\}$，所有符号的虚部为零；转移函数直接模拟$M_{\det}$。由于所有符号虚部为零，永不触发分支，生成元计数器始终为零，且没有生成元被引入和激活。故对所有$x$有$\kappa_M(x)=0$。由$\kappa(L)$的定义，$\kappa(L)\le 0$，因此$\kappa(L)=0$。
	\end{proof}
	
	我们建立了第一个关键证据：\textbf{所有$\mathbf{P}$类语言的本质维度为零，即有界的。}这对应于``高效可解的问题无需本质的代数扩张''。P类问题的验证器可以在基域内完成所有计算，无需触发任何非确定性分支，因此不消耗任何代数维度资源，也不需要任何代数对偶结构。
	
	\section{$\kappa(L)$的障碍无关性}
	
	本节系统论证本质维度$\kappa(L)$不受相对化、自然证明和代数化三大障碍的约束，为基于$\kappa(L)$的分离论证提供方法论保障。三大障碍的核心均在于传统图灵机对谕示、伪随机函数和代数闭包的透明性，而CBTM/IVM框架通过投影约束、定义域限制和阈值离散性三个机制，在操作语义层面填补了经典模型的语义空洞，使得这些障碍的前提被逐一瓦解。
	
	\subsection{投影约束与相对化障碍}
	
	相对化障碍的核心在于传统图灵机对谕示响应不加任何过滤，谕示返回的比特可以直接写入纸带，从而影响计算的任何部分。在CBTM/IVM框架中，投影约束公理（见论文\cite{models2026}）强制写入符号的虚部完全由当前状态和读取符号的虚部通过布尔函数决定，而不能由外部直接赋值。
	
	\begin{definition}[谕示响应编码]\label{def:oracle-encoding}
		对任意谕示$O$及查询$q$，响应$\tau_O$编码为$\tau_O = \text{value} + 0\cdot\alpha$，即$\mathfrak{Im}(\tau_O)=0$。谕示响应被\textbf{锁死在实部}，虚部恒为零。
	\end{definition}
	
	谕示查询步骤因此不会引入新的代数生成元，不会改变已有生成元之间的线性关系，无法影响虚带的演化。在计算的核心周围形成了一层``无菌隔离区''：外部世界可以提供数据，但无法带入新的代数维度。谕示只能影响实部（确定性数据部分），而无法触及虚部中的不可公度性和对偶性——这正是非确定性的本质所在。
	
	从更根本的层面看，经典NP定义所依赖的“猜测虚部控制信息并拷贝到实部”这一不可能操作，在相对化世界中可以被谕示免费完成——谕示通过猜测获得见证并直接将其作为数据返回给机器，从而填补了经典模型操作不自足所留下的语义空洞。这正是相对化世界中出现 $\mathbf{P}^A = \mathbf{NP}^A$ 假象的深层根源。投影约束机制通过将谕示响应强制锁定在实部，从根本上杜绝了谕示承担这一不可能操作的可能性，使 $\kappa(L)$ 在任何谕示环境下保持稳定。
	
	\begin{theorem}[$\kappa(L)$的谕示不变性]\label{thm:oracle-invariance}
		设$L$为语言，$O$为任意谕示。在投影约束和谕示响应编码规则下，有$\kappa(L^O) = \kappa(L)$。
	\end{theorem}
	\begin{proof}
		设$M$为在代数域$K = \mathbb{Q}(\sqrt{p_1}, \dots, \sqrt{p_k})$上展开计算的IVM。考虑$M^O$。谕示查询步骤的处理如下：机器暂停正常计算，编码查询并发送给$O$；谕示返回比特$b \in \{0,1\}$，被编码为符号$\tau_O = b + 0 \cdot \alpha$（虚部恒为零）；$\tau_O$被写入当前工作带单元，该操作只影响实带，虚带保持不动。
		
		由于$\mathfrak{Im}(\tau_O)=0$，投影约束保证后续任何写入符号的虚部仅由当前状态和读取符号的虚部（为零）决定，且状态本身不携带生成元信息。因此，整个计算过程中虚带的演化与无谕示机器$M$完全一致：没有因谕示响应而引入的新生成元，已有生成元的激活状态不变。因而对所有输入$x$，$\kappa_{M^O}(x) = \kappa_M(x)$。由于$M$任意且$\kappa(L)$定义为所有验证机的下确界，故$\kappa(L^O) = \kappa(L)$。
	\end{proof}
	
	投影约束机制确保$\kappa(L)$的值不随谕示的引入而改变，使相对化障碍的前提不再成立。$\kappa(L)$所揭示的$\mathbf{P}$与$\mathbf{NP}$的代数资源鸿沟在所有相对化世界中依然潜在存在。
	
	\subsection{定义域限制与自然证明障碍}
	
	自然证明障碍要求下界证明必须基于对大多数布尔函数成立的``大尺度''组合性质。我们的核心观察是，$\kappa(L)$\textbf{根本上不满足``大尺度''条件}，并且这并非由于技术困难，而是因为它对\textbf{绝大多数布尔函数逻辑上无定义}。
	
	\begin{theorem}[$\kappa(L)$的定义域限制]\label{thm:domain-restriction}
		$\kappa(L)$有定义当且仅当$L \in \mathbf{NP}$。特别地，若$L \notin \mathbf{NP}$，则$\mathcal{V}_L = \emptyset$且$\kappa(L)$无定义。
	\end{theorem}
	
	考虑全体$n$元布尔函数集合$\mathcal{F}_n$。每个$f\in\mathcal{F}_n$对应一个语言$L_f \subseteq \{0,1\}^n$。
	
	\begin{lemma}[$\mathbf{NP}$语言的稀疏性]\label{lem:np-sparsity}
		存在多项式$p$使得$|\mathbf{NP}_n| \le 2^{p(n)}$。因此，
		\[
		\Pr_{f\sim\mathcal{F}_n}[L_f \in \mathbf{NP}] \le \frac{2^{p(n)}}{2^{2^n}} \to 0 \quad (n\to\infty).
		\]
	\end{lemma}
	\begin{proof}
		每个$\mathbf{NP}$语言可由一台多项式时间验证机描述，而验证机的描述长度是多项式的，因此数目至多为$2^{\mathrm{poly}(n)}$。
	\end{proof}
	
	\begin{corollary}[$\kappa$对随机函数几乎处处无定义]\label{cor:kappa-undefined}
		对随机布尔函数$f$，$\kappa(L_f)$有定义的概率趋于$0$。
	\end{corollary}
	
	因此，性质$\mathcal{P}_{\text{high}}(L)$定义为``$\kappa(L)=\infty$''不满足``大尺度''条件。$\mathcal{P}_{\text{high}}$不是自然性质。
	
	\begin{remark}[代数根源]
		自然性质基于布尔函数的组合属性，完全处于标准图灵机的表达能力范围内。而$\kappa(L)$依赖于生成元的线性无关性，这是一个超出布尔函数组合范畴的代数概念。这种``范畴错位''正是自然证明障碍对$\kappa(L)$失效的根本原因。此外，判定$\kappa(L)$的值（例如是否大于$0$）本身等价于$\mathbf{P}$ vs.\ $\mathbf{NP}$问题，故极可能不具有$2^{O(n)}$构造性，这进一步排除了落入自然证明框架的可能性。
	\end{remark}
	
	定义域限制从根本上规避了自然证明障碍。该论证是结构性的、无条件的，不依赖于任何伪随机性假设。
	
	\subsection{阈值离散性、伽罗瓦敏感性与代数化障碍}
	
在讨论代数化障碍的细节之前，有必要指出其与相对化障碍的深层关系。 从机制上看，代数化框架（Aaronson \& Wigderson）是将传统谕示替换为低次多项式扩展，其核心运载工具仍然是外部谕示。因此，代数化障碍本质上继承自相对化障碍：任何能够切断机器与外部谕示之间‘语义空洞传递’的机制，若对普通谕示有效，则对代数化谕示同样构成屏蔽。本文所确立的投影约束公理，恰恰提供了这样一种机制——它强制一切外部响应（无论是否经过多项式扩展）的虚部恒为零。这意味着，破除相对化障碍的那个公理，在代数化的框架下仍然牢不可破。
		
	2009年，Aaronson和Wigderson~\cite{aaronson2009}将谕示扩展为低次多项式，证明任何对谕示扩展保持有效的证明技术（``代数化''证明）无法分离$\mathbf{P}$和$\mathbf{NP}$。然而，代数化扩展的效力范围仅限于软化谕示的输入输出关系，无法改变机器内部与谕示调用无关的基本操作。
	
	在CBTM/IVM框架中，分支触发由阈值判断$\mathfrak{Im}(\tau) > \epsilon$（$\epsilon=0.5$）决定。该判断具有本质的\textbf{完全离散性}：
	\begin{enumerate}
		\item \textbf{输入离散}：$\mathfrak{Im}(\tau)$只能取$0$或$1$，由机器自身写入的符号决定，与谕示响应无关。
		\item \textbf{比较离散}：阈值$\epsilon=0.5$是固定有理数，比较运算不存在``部分成立''的中间状态。
		\item \textbf{结果离散}：判断结果只有``分支''或``不分支''两种，直接决定计算树的拓扑结构。
	\end{enumerate}
	
	阈值判断检测的是``虚部是否非零''，本质上是在检测\textbf{不可公度元是否被引入}。这是一个代数判断，而非组合判断。代数化谕示所能回答的所有问题，归根结底都是关于\textbf{多项式等式}的问题。$\sqrt{p_i}$与$-\sqrt{p_i}$满足完全相同的代数方程$x^2 - p_i = 0$，任何代数查询都无法区分它们。然而，阈值判断却对它们做出不同响应：$+1 > 0.5$触发分支，$-1 < 0.5$不触发。这意味着，非确定性分支的触发依赖于符号的正负号——即对偶生成元的选择（$+\sqrt{p_i}$ 还是 $-\sqrt{p_i}$）——而非多项式等式。代数化谕示对此天然``失明''。
	
	如果$\mathbf{P} = \mathbf{NP}$，那么非确定性分支应该可以被确定性模拟消除。但消除的前提是，确定性算法能够区分``触发分支''与``不触发分支''的代数条件。这种条件依赖于符号信息，而任何仅依赖多项式等式的代数化框架都无法感知这种信息。因此，代数化框架下的任何构造，都无法在CBTM/IVM中强制$\mathbf{P} = \mathbf{NP}$。$\mathbf{P}$与$\mathbf{NP}$的差异恰好寓于代数化框架的``盲点''之中——这一盲点正是对偶性（生成元的正负号差异），它在经典NP定义中被存在量词消除，在代数化框架中无法被感知，而在IVM中被显式保留并成为分离的关键。
	
	\begin{definition}[伽罗瓦敏感性]\label{def:galois-sensitivity}
		设$M$为CBTM/RBTM。称$M$在输入$x$上\textbf{对伽罗瓦自同构敏感}，如果存在接受路径$\pi$，使得其对偶路径$\sigma(\pi)$（翻转所有非确定性选择）不是接受路径。否则称$M$\textbf{不敏感}。
	\end{definition}
	
	\begin{lemma}[敏感性的谕示不变性]\label{lem:sensitivity-invariance}
		设$M$为一台CBTM/RBTM，$O$为任意谕示，$\tilde{O}$为其任意低次多项式扩展。则$M$在输入$x$上的伽罗瓦敏感性在$O$和$\tilde{O}$下完全相同。
	\end{lemma}
	\begin{proof}
		路径切换仅改变分支选择序列$\mathbf{d} \mapsto \mathbf{d} \oplus \mathbf{1}_m$。该操作不涉及谕示查询，不改变实部内容，也不依赖谕示返回值。因此，$\pi$接受当且仅当$\sigma(\pi)$接受，这一事实与谕示及其扩展无关。
	\end{proof}
	
	\begin{theorem}[P类的不敏感性]\label{thm:p-insensitivity}
		若$L \in \mathsf{P}$，则存在判定$L$的机器在所有输入上完全不敏感。
	\end{theorem}
	\begin{proof}
		$L \in \mathsf{P}$意味着存在$\mathrm{CBTM}|_0$判定$L$，其所有符号虚部恒为零，计算树退化为单一路径，无非确定性选择点。路径切换是恒等操作。
	\end{proof}
	
	\subsection{谕示无关证明}
	
	\begin{theorem}[敏感性证明的非代数化性]\label{thm:non-algebrizing}
		基于伽罗瓦自同构敏感性的$\mathsf{P} \neq \mathsf{NP}$证明是\textbf{非代数化}的，不受Aaronson-Wigderson代数化障碍的约束。
	\end{theorem}
	\begin{proof}
		代数化障碍的成立依赖于一个隐含前提：所有被考虑的证明技术都在标准图灵机模型下运作，且非确定性的力量可以通过谕示来操纵，因为谕示响应可以影响计算的任何部分。CBTM/IVM框架通过以下机制打破了这一前提：
		\begin{enumerate}
			\item \textbf{投影约束}将谕示响应强制编码为虚部为零的符号，使谕示只能影响实部计算；
			\item \textbf{阈值判断的完全离散性}确保分支触发完全由机器内部符号的虚部决定，且永不因谕示查询而改变；
			\item \textbf{谕示无关性}（引理\ref{lem:sensitivity-invariance}）进而保证，伽罗瓦敏感性在任何谕示环境下都保持不变。
		\end{enumerate}
		
		因此，CBTM/IVM模型中的非确定性力量被重新定位在谕示无法触及的代数结构（虚部空间与域扩张）中。任何在标准图灵机模型下使$\mathsf{P} = \mathsf{NP}$的谕示构造（包括代数化谕示），都无法破坏CBTM/IVM中的阈值判断与敏感性。
		
		然而，这并不使该证明成为``代数化''的。关键区别在于：代数化证明的特征是，其推理步骤主动利用谕示的扩展性质，因而其有效性依赖于对谕示代数化的适应。我们的证明恰恰相反，其核心推理步骤（阈值判断、分支触发、路径切换）完全绕开了谕示，对谕示的扩展既不依赖也不适应，仅仅是因为谕示扩展无法触及这些步骤，才使得结论在扩展下被动地得以保持。这是一个\textbf{谕示无关证明}，而非代数化证明。
	\end{proof}
	
	\subsection{障碍无关性总结}
	
	三大障碍的共同根源在于经典语言无法表达不可公度性和对偶性。相对化障碍的根源在于谕示只能返回基域内的布尔值，无法触及虚部中的不可公度性；自然证明障碍的根源在于自然性质基于组合属性，而本质维度是代数概念，超出了组合范畴；代数化障碍的根源除了相对化障碍中的问题以外，还在于代数化谕示只能回答多项式等式问题，无法感知阈值判断所依赖的符号信息（生成元的正负号，即对偶性）。CBTM/IVM框架通过内建不可公度性并显式保留对偶性，在操作语义层面填补了经典模型的语义空洞，使这些障碍的前提被逐一瓦解。
	
	$\kappa(L)$被严格证明为\textbf{障碍无关的不变量}，能够区分$\mathbf{P}$与$\mathbf{NP}$而不受三大障碍干扰。这为$\mathbf{P} \neq \mathbf{NP}$的证明扫清了所有已知的方法论障碍。
	
	
	综上，代数化谕示之所以无法扰动 IVM 的分支行为，根源在于它像所有谕示一样，被投影约束公理严格限制在实部。代数化扩展并未赋予它绕过虚部禁区的额外能力。因此，克服代数化障碍不需要全新的技术，只需要严格贯彻从相对化障碍中就已确立的同一结构法则。”
	
	\section{NP完全问题的维度下界}
	
	本节给出NP完全问题本质维度下界的两个独立论证。第一个论证基于信息论下界与分支-激活引理，从组合角度证明任何正确验证器都必须经历至少$n$次非确定性分支，从而激活至少$n$个线性无关的生成元。第二个论证基于验证敏感性定理，从代数角度证明生成元的激活是验证正确性的必要条件，不可消除。两个论证互补，共同加固下界的普适性和不可绕过性。
	
	\subsection{子集和的信息论下界}
	
	设 $M$ 为任意正确验证子集和问题的 IVM。考虑 $n$ 元素实例 $(S, t)$ 的标准编码，其中 $S = \{s_1, \dots, s_n\}$ 由互异正整数组成。每个元素 $s_i$ 被编码为 $\tau_i = s_i + 1 \cdot \sqrt{p_i}$（$p_i$ 为第 $i$ 个素数）。在此编码下，$\tau_i$ 的虚部为 $1$，表明当验证器读取该符号时将触发非确定性分支。关键之处在于，虚部标记仅依赖于 $S$（通过固定的素数序列），而与目标值 $t$ 完全无关。根据投影约束公理，虚部的演化完全由当前状态和纸带上已有的虚部决定，绝不会受 $t$ 的影响。因此，对于固定的集合 $S$，机器 $M$ 在任何输入 $(S, t)$ 上的计算树都具有完全相同的分支拓扑；该树由 $S$ 唯一确定，不随 $t$ 的改变而改变。
	
	\begin{lemma}[路径数下界]\label{lem:path-count}
		设 $S$ 为 $n$ 个互异正整数的集合，$M$ 为任意正确验证子集和的 IVM。在所有实例 $(S, \cdot)$ 共享的计算树上，不同计算路径的条数至少为 $2^n$。
	\end{lemma}
	\begin{proof}
		对每个子集 $T \subseteq S$，定义 $t_T = \sum_{s_i \in T} s_i$。因为 $T$ 本身是见证，实例 $(S, t_T)$ 是 YES 实例，$M$ 必须接受它。因此在共享的计算树中至少存在一条接受路径；我们选取这样一条路径并记为 $\pi_T$。
		
		假设共享树中不同路径的条数少于 $2^n$。由于共有 $2^n$ 个不同的子集 $T$，由鸽巢原理，必有两个不同的子集 $T_1 \neq T_2$ 被映射到同一条接受路径，即 $\pi_{T_1} = \pi_{T_2}$。这两条路径对应于相同的分支选择序列 $\mathbf{d} \in \{0,1\}^m$。
		
		因为 $M$ 在分支选择固定后是确定性的，沿着 $\mathbf{d}$ 的计算会产生一个确定的子集和 $s(\mathbf{d})$。在实例 $(S, t_{T_1})$ 中，沿 $\mathbf{d}$ 运行导致接受，故 $s(\mathbf{d}) = t_{T_1}$。在实例 $(S, t_{T_2})$ 中，同一条路径同样导致接受，故 $s(\mathbf{d}) = t_{T_2}$。于是 $t_{T_1} = t_{T_2}$。但 $S$ 中元素互异且为正整数，不同子集对应不同和值，这与 $T_1 \neq T_2$ 矛盾。
		
		因此，共享树必须包含至少 $2^n$ 条不同路径。
	\end{proof}
	
	\begin{lemma}[分支次数下界]\label{lem:branch-count}
		设 $M$ 为任意正确验证子集和的 IVM。对于任意 $n$ 元素实例 $(S, t)$，$M$ 的计算树中至少包含 $n$ 个非确定性分支。
	\end{lemma}
	\begin{proof}
		由引理 \ref{lem:path-count} 知，该树至少有 $2^n$ 条不同路径。计算树是二叉的（每个分支恰好产生两个后继），而一棵含有 $m$ 个内部非确定性节点的二叉树至多有 $2^m$ 个叶节点（不同路径）。因此 $2^m \ge 2^n$，得 $m \ge n$。
	\end{proof}
	
	
	\subsection{分支-激活引理与本质维度下界}
	
	\begin{lemma}[分支-激活引理，重述]\label{lem:branch-activate-2}
		设$M$为任意IVM。若在某个计算步骤中发生非确定性分支，则该分支必然激活恰好一个生成元。具体而言，在分支触发的瞬间，IVM自动分配一个新的素数平方根地址$\sqrt{p_i}$（同时对偶地引入$-\sqrt{p_i}$），并在激活路径上记录值$1$，对应的生成元被激活。这一记账是IVM的定义性行为，独立于底层转移函数的设计，无法被任何算法绕过。
	\end{lemma}
	
	\begin{theorem}[子集和的本质维度下界]\label{thm:subset-sum-lower}
		$\kappa(\mathrm{SubsetSum}) = \Omega(n)$。具体地，对任意正确的IVM验证器$M$以及任意$n$元素实例，$\kappa_M(x_n) \ge n$。
	\end{theorem}
	\begin{proof}
		设$M$为任意正确验证子集和问题的IVM。对任意$n$元素实例$(S, t)$，由引理\ref{lem:branch-count}，$M$在验证过程中至少发生$n$次非确定性分支。由引理\ref{lem:branch-activate-2}，每次分支激活恰好一个生成元。因此至少$n$个生成元被激活，这些生成元为互异素数平方根$\{\sqrt{p_1}, \dots, \sqrt{p_n}\}$。由素数平方根线性无关定理（引理\ref{lem:commensurability}），这$n$个生成元在$\mathbb{Q}$上线性无关，故它们张成的空间维数至少为$n$，即$\kappa_M(S, t) \ge n$。
		
		由于$M$是任意正确验证器，取所有$M \in \mathcal{V}_{L_{\mathrm{SubsetSum}}}$的下确界，得
		\[
		\kappa(L_{\mathrm{SubsetSum}}) = \inf_{M \in \mathcal{V}_L} \sup_{(S,t)} \kappa_M(S, t) \ge n = \Omega(n).
		\]
	\end{proof}
	
	该定理的证明逻辑链极为清晰：信息论下界强制了分支次数（$\ge n$），分支-激活引理将分支次数转化为生成元激活数量（恰好$n$个），素数平方根线性无关定理将激活数量转化为代数维数（$\ge n$维）。三个步骤各自独立，又环环相扣，构成了一个不可绕过的下界论证。
	
	\subsection{验证敏感性定理}
	
	上述信息论下界论证确立了分支次数和生成元数量足够多（$\ge n$）。验证敏感性定理从互补角度证明：这些生成元的激活不是冗余的，而是验证正确性的必要条件——单点翻转必然改变某些实例的判定结果。两者共同确保了子集和下界的普适性和不可绕过性：前者堵死了``分支不够多''的可能，后者堵死了``分支虽多但都冗余''的可能。
	
	\begin{theorem}[验证敏感性定理]\label{thm:sensitivity}
		设$M$为任意正确验证子集和问题的IVM。则存在一个输入实例以及该实例计算树中的一个非确定性选择点，使得单独翻转该点的选择能够改变计算结果（接受变为拒绝，或拒绝变为接受）。特别地，不存在完全不敏感的验证器能够判定子集和问题。
	\end{theorem}
	\begin{proof}
		构造专门实例。取$S = \{s_1,\dots,s_n\}$，所有元素互异且为正整数。对任意固定的$i \in \{1, \dots, n\}$，令目标值$t = s_i$。此实例的唯一解为$\{s_i\}$。
		
		$M$必须接受该实例。设$\pi$为接受路径，其分支选择序列为$\mathbf{d} = (d_1, \dots, d_m)$。由引理\ref{lem:branch-count}，$m \ge n$。在第$i$个非确定性分支处，$\pi$做出了一个确定的选择$d_i \in \{0, 1\}$。施加单点翻转$\rho_i$：将$d_i$替换为$1 - d_i$，保持其他选择不变，得到路径$\pi^{\oplus i}$。
		
		$\pi^{\oplus i}$对应一个在元素$s_i$上选择被反转后的子集。由于唯一解为$\{s_i\}$，$\pi^{\oplus i}$对应的子集必然不等于$\{s_i\}$，其子集和$\neq t = s_i$。因此$\pi^{\oplus i}$不是解，$M$对$\pi^{\oplus i}$必须给出``拒绝''判定。故$M$在专门实例$(S, s_i)$上对第$i$个选择点单点敏感。
		
		由于$i$是任意的，每个非确定性选择点都有专门实例使其敏感。故$M$不可能完全不敏感——任何正确验证器必然在某些输入上对某些选择点单点敏感。翻转必然改变结果，意味着生成元的激活状态确实参与并影响了接受/拒绝的判定，因此生成元的激活是验证正确性的必要条件，不可消除。
	\end{proof}
	
	验证敏感性定理揭示了$\mathbf{NP}$完全问题的验证器必然具有``牵一发而动全身''的结构——每个非确定性选择点在某个输入下都是决定性的。这与$\mathbf{P}$类语言形成鲜明对比：$\mathbf{P}$类语言可以完全不依赖非确定性选择（定理\ref{thm:p-insensitivity}），因而对任何选择点的翻转都不敏感。正是这种结构性差异将两类语言分离开来。从对偶性的角度看，单点敏感性正是对偶性在单个选择点上的体现：如果每个选择点的对偶翻转都不影响结果，那么对偶性就是冗余的，非确定性可以被消除；反之，如果存在某个选择点其对偶翻转必然改变结果，那么对偶性就是本质的，非确定性无法被消除。
	
	\subsection{维度保持归约与所有NP完全问题的无界维度}
	
	维度爆炸并非子集和的孤立现象。多项式时间归约是连接$\mathbf{NP}$完全问题的基本工具。以下定理证明，归约不会压缩本质维度，因此子集和的$\Omega(n)$下界可以传递至所有$\mathbf{NP}$完全问题。
	
	\begin{theorem}[维度保持归约]\label{thm:dimension-preserving}
		若$A \le_p B$，则$\kappa(B) \ge \kappa(A)$。
	\end{theorem}
	\begin{proof}
		任取验证$B$的IVM $M_B \in \mathcal{V}_B$。由多项式时间归约的定义，存在确定性图灵机$M_f$在多项式时间内计算归约函数$f$。由定理\ref{thm:p-zero}，$M_f$可转化为零维IVM，其在计算过程中永不触发分支，虚部恒为零。构造验证$A$的IVM $M_A$如下：对输入$x$，
		\begin{enumerate}
			\item 计算$y = f(x)$，其中$f$由零维IVM实现（不引入生成元）。
			\item 模拟$M_B$在输入$y$上的验证过程。
		\end{enumerate}
		由于$f$的计算不引入生成元，$M_A(x)$的虚部空间与$M_B(y)$的虚部空间完全相同。因此$\kappa_{M_A}(x) = \kappa_{M_B}(f(x))$。
		
		需要验证$M_A$在多项式时间内运行。因为$A \le_p B$，存在多项式$p$使得$|y| = |f(x)| \le p(|x|)$。$M_B$的运行时间为$T_B(|y|)$，为多项式，故$M_A$模拟$M_B$的总步数至多为$T_B(p(|x|))$，仍是$|x|$的多项式。加上计算$f(x)$的多项式时间，$M_A$整体在多项式时间内运行。
		
		对$x$取上确界，得
		\[
		\sup_x \kappa_{M_A}(x) = \sup_x \kappa_{M_B}(f(x)) \le \sup_y \kappa_{M_B}(y).
		\]
		右边对$M_B$取下确界即得$\kappa(B)$，而左边至少为$\kappa(A)$。故$\kappa(A) \le \kappa(B)$。
	\end{proof}
	
	\begin{corollary}[所有NP完全问题具有无界维度]\label{cor:np-complete-unbounded}
		对任意$\mathbf{NP}$完全语言$L$，$\kappa(L) = \Omega(n)$。
	\end{corollary}
	\begin{proof}
		子集和问题$L_{\mathrm{SubsetSum}}$是$\mathbf{NP}$完全的，故$L_{\mathrm{SubsetSum}} \le_p L$。由定理\ref{thm:dimension-preserving}，$\kappa(L) \ge \kappa(L_{\mathrm{SubsetSum}}) = \Omega(n)$。
	\end{proof}
	
	维度爆炸并非子集和的孤立现象，而是所有$\mathbf{NP}$完全问题的共同特征——\textbf{每个$\mathbf{NP}$完全问题都需要无界的代数维度资源。}这构成了$\mathbf{P}$与$\mathbf{NP}$之间的一个结构性鸿沟：$\mathbf{P}$类语言的本质维度有界（$\kappa=0$），而$\mathbf{NP}$完全语言的本质维度无界（$\kappa=\Omega(n)$）。
	
	\section{$\mathbf{P} \neq \mathbf{NP}$的证明}
	
	\subsection{维度分离原理}
	
	综合第3节和第5节的结论，我们得到了一个精确的维度分离图景：
	\begin{itemize}
		\item 对所有$L \in \mathbf{P}$，$\kappa(L) = 0$（定理\ref{thm:p-zero}）。
		\item 对所有$\mathbf{NP}$完全语言$L$，$\kappa(L) = \Omega(n)$（推论\ref{cor:np-complete-unbounded}）。
	\end{itemize}
	
	线性相关（零维）与线性无关（$\Omega(n)$维）是互斥的代数结构。一个代数空间不可能同时是零维和$\Omega(n)$维——前者意味着没有任何生成元被激活，后者意味着至少$\Omega(n)$个线性无关的生成元被激活。这一互斥性构成了反证法的逻辑基础。
	
	\subsection{反证法}
	
	\begin{theorem}[$\mathbf{P} \neq \mathbf{NP}$]\label{thm:pnp}
		$\mathbf{P} \neq \mathbf{NP}$。
	\end{theorem}
	\begin{proof}
		假设$\mathbf{P} = \mathbf{NP}$。则子集和问题$L_{\mathrm{SubsetSum}} \in \mathbf{NP}$蕴含$L_{\mathrm{SubsetSum}} \in \mathbf{P}$。由定理\ref{thm:p-zero}，$\kappa(L_{\mathrm{SubsetSum}}) = 0$。但由定理\ref{thm:subset-sum-lower}，$\kappa(L_{\mathrm{SubsetSum}}) = \Omega(n)$。对于充分大的$n$，$0$不可能等于渐近增长的函数，矛盾。故假设不成立，$\mathbf{P} \neq \mathbf{NP}$。
	\end{proof}
	
	该证明极简而有力。它不需要构造具体的语言来见证分离，不需要复杂的对角线论证，也不需要精细的电路下界分析。它仅依赖于一个根本性的观察：$\mathbf{P}$类语言可以在基域内完成所有计算（零维），而$\mathbf{NP}$完全语言的验证本质上需要不断扩张的代数维度（$\Omega(n)$维）。这两者之间的鸿沟不是量的差异，而是质的对立——线性相关与线性无关在代数上是互斥的。
	
	\subsection{障碍无关性的重申}
	
	此证明完全建立在障碍无关的不变量$\kappa(L)$之上。第4节已详细论证了$\kappa(L)$在相对化下稳定、仅对$\mathbf{NP}$语言有定义且具有伽罗瓦敏感性。具体而言：
	
	\begin{itemize}
		\item \textbf{相对化障碍}：从CBTM的视角来看，相对化障碍的本质根源在于：谕示可以被暗中赋予填补语义空洞的超能力——替机器完成“猜测虚部控制信息并拷贝到实部”这一经典模型自身不可能执行的操作。$\mathbf{P}^A = \mathbf{NP}^A$ 的假象之所以能在相对化世界中出现，正是因为谕示免费承担了经典NP定义所依赖但经典操作语义不可实现的那一拷贝操作。投影约束机制通过强制谕示响应编码为虚部为零的符号，将谕示限制在实部，确保 $\kappa(L)$ 在任何谕示环境下保持不变。谕示只能影响实部（确定性数据部分），而无法触及虚部中的不可公度性和对偶性。因此基于 $\kappa(L)$ 的分离结论在所有相对化世界中依然成立。
		
		\item \textbf{自然证明障碍}：$\kappa(L)$依赖于生成元的线性无关性，这是一个超出布尔函数组合范畴的代数概念。$\kappa(L)$仅对$\mathbf{NP}$语言有定义，对随机布尔函数几乎处处无定义，从而不满足自然证明的``大尺度''条件。
		
		\item \textbf{代数化障碍}：代数化谕示的核心操作是将谕示替换为低次多项式扩展。这种扩展使得谕示能够回答关于多项式等式的任意问题，但它被限制在一种非常特定的数学形式中：它只能返回多项式等式的值，无法感知符号的正负号。对于代数化谕示来说，$+1$和$-1$是完全不可区分的——它们满足完全相同的代数方程 $x^2-1=0$。因此，无论你如何构造一个能让 $\mathbf{P}=\mathbf{NP}$ 的代数化谕示，它都无法触及CBTM/IVM中 $\mathbf{P}$ 与 $\mathbf{NP}$ 之间的本质差异。因为这个差异恰好寓于代数化谕示的感知盲区之中——那个盲区正是对偶性。
		
		在CBTM/IVM中，非确定性分支的触发由阈值判断 $\mathfrak{Im}(\tau) > 0.5$ 决定，其核心在于区分虚部的正负号——即对偶生成元的选择。代数化谕示所能回答的所有问题都是关于多项式等式的，无法感知符号信息。因此代数化谕示对触发分支与否的关键信息天然``失明''。基于敏感性的证明是谕示无关证明，不受代数化障碍约束。
	\end{itemize}
	
	我们以数学的确定性得出结论：
	\[
	\boxed{\mathbf{P} \neq \mathbf{NP}}.
	\]
	此结论不仅解决了一个具体的数学问题，更为理解计算的本质提供了深刻的代数几何洞察。整个证明在ZFC公理体系内完成，不依赖于任何未经证实的数学猜想。证明的核心——信息论下界、分支-激活引理、素数平方根线性无关定理——各自都是经典数学中已被严格确立的结果，它们的组合构成了一个不可动摇的论证链条。
	
	\section{讨论：重访经典框架}
	
	\subsection{为什么经典框架无法证明$\mathbf{P} \neq \mathbf{NP}$}
	
	本文的证明建立在CBTM/IVM框架之上，该框架的核心特征是内建了不可公度性作为操作语义的内在属性，并显式保留了对偶性，同时从根本上消除了拷贝操作的必要性。一个自然的问题是：既然参数化等价定理保证了CBTM/IVM与经典图灵机在识别语言类上的外延等价（$\mathbf{P}_{cb}=\mathbf{P}$，$\mathbf{NP}_{cb}=\mathbf{NP}$），为什么这一证明不能被视为经典框架内的证明？
	
	答案在于\textbf{外延等价不等于内涵等价}。CBTM/IVM与经典模型识别相同的语言类，但两者的操作语义有本质差异。我们分几个层次来阐述经典框架无法复制这一证明的根本原因。
	
	\textbf{第一：经典NP定义在定义过程中已丢失了非确定性的核心语义。} 在进入编码和维度消耗的讨论之前，我们必须首先注意到一个更为根本的事实：经典NP定义本身，在将见证$c$写在实部纸带上的那一刻，就已经丢失了非确定性选择的两个核心语义属性——\textbf{不可公度性}和\textbf{对偶性}。不可公度性在函数指针被扭曲为数据时丢失——$c$的各个比特之间不再保留任何代数独立性关系，经典模型无法从数学上区分``真正独立的非确定性选择''与``仅在语法标签上不同的选项''。对偶性在编码为比特串并被存在量词包裹时丢失——$+\sqrt{p_i}$与$-\sqrt{p_i}$的对偶关系在比特$0$和$1$中没有对应物，未被选择的路径的所有痕迹都被抹去。因此，经典框架中的``非确定性''从一开始就是一个语义残缺的概念。任何试图使用这个残缺概念来证明P与NP分离的努力，都面临一个原则性的困难：论证所依赖的核心概念本身已在定义中被阉割。
	
	\textbf{第二，编码的非中立性。} 经典模型可以通过非平凡编码来``表示''不可公度性（在弱意义上判定它），但无法将其``内建''为操作语义的内在属性（在强意义上锚定它）。编码的本质是外延指代，而非内涵内建。正如论文\cite{models2026}所论证的，语法模拟的成功并不意味着语义的保持——当面对连续量或代数不可公度元素时，编码会丢失关键的数学信息。特别是，经典框架通过将见证$c$从虚部拷贝到实部（这一拷贝本身在经典操作语义中不可实现），完成了从函数指针到数据的身份扭曲，同时丢失了不可公度性；并用存在量词消除了对偶性——这些概念错误使得任何在经典框架内通过编码来``模拟''IVM的维度计数的尝试，都面临一个根本性的困难：编码后的操作依赖于外部语义赋值，而非由机器的操作语义所内建。
	
	\textbf{第三，维度消耗的强制性。} IVM中的分支-激活引理所确立的维度消耗是操作语义的强制性记账——正如经典图灵机的步数统计无法被转移函数``关闭''，IVM的生成元激活记录同样无法被绕过。这种强制性在经典模型中缺乏对应物，因为经典模型的操作语义根本不包含``生成元激活''这一概念。在经典模型中，你可以定义``步数''，但你无法定义``生成元激活''——因为经典模型的操作语义无法区分``这个符号涉及不可公度元''和``这个符号不涉及不可公度元''（这正是不可公度性元定理的结论），也无法追踪对偶性的保留与否。
	
	因此，本文的证明不能被视为经典框架内的证明。它使用了经典模型所缺乏的概念资源——不可公度性作为操作语义的内在属性、对偶性的显式保留、本质维度$\kappa(L)$——并且修正了经典NP定义的操作不自足和概念错误。在这个意义上，经典P vs.\ NP问题在经典操作语义内部是``不可证''的——不是因为该命题在绝对意义上是独立的，而是因为经典语言缺乏表达分离所需的核心概念，且经典NP定义的概念错误使得这些核心概念在经典框架中不可定义。
	
	\subsection{关于经典P vs.\ NP问题的定位}
	
	基于上述分析，我们可以精确地定位本文的证明与经典P vs.\ NP问题的关系。
	
	\textbf{第一：经典NP定义在操作上不自足，在语义上残缺。} 该定义要求了一个在经典操作语义中不可实现的操作（虚部到实部拷贝），且在不可避免的编码过程中丢失了非确定性的两个核心语义属性——不可公度性（第一重错误的后果）和对偶性（第二重错误的后果）。这一操作不自足和语义残缺是经典框架所有困难的最终根源。CBTM/IVM框架通过将虚部内建为操作语义、消除拷贝操作、保留不可公度性和对偶性，从根本上修正了这些缺陷。
	
	\textbf{第二，$\mathbf{P} \neq \mathbf{NP}$在集合论意义下为真。} 参数化等价定理保证了$\mathbf{P}_{cb} = \mathbf{P}$和$\mathbf{NP}_{cb} = \mathbf{NP}$。由于我们在CBTM/IVM框架中严格证明了$\mathbf{P}_{cb} \neq \mathbf{NP}_{cb}$，因此作为集合论命题，$\mathbf{P} \neq \mathbf{NP}$为真。该真值不依赖于我们使用什么计算模型来证明它。P和NP作为语言集合，它们要么相等，要么不相等——这是二值逻辑下的客观事实。我们的工作确立了这个事实的真值为假。
	
	\textbf{第三，该证明在ZFC内完成。} 本文的整个证明所涉及的证明——不可公度性元定理（论文\cite{models2026}）、CBTM/IVM框架的构造、信息论下界论证、素数平方根线性无关定理、反证法——全部在ZFC公理体系内完成。我们没有使用任何超出ZFC的假设（如大基数公理、选择公理的强化版本等）。$\mathbf{P} \neq \mathbf{NP}$在ZFC内可证，不是独立于ZFC的命题。
	
	\textbf{第四，该证明无法在经典操作语义内部被复制。} 证明所需的核心概念——不可公度性作为操作语义的内在属性、对偶性的显式保留、本质维度$\kappa(L)$——在经典操作语义中不可定义。经典NP定义的概念错误（见证本体论错位导致不可公度性丢失、对偶性消除、操作不自足）使得这些概念在经典框架中根本没有对应物。因此，任何局限于经典语言内的证明都无法完成分离。这不是说该问题在经典框架内没有真值，而是说经典框架缺乏证明它所需的概念资源。
	
	\textbf{第五，该问题在经典框架内被``消解''。} ``消解''一词在此有精确的含义：它不意味着取消问题的真值，而是揭示了原初问题的隐含预设（经典模型能正面定义非确定性独立性）为假，并诊断了经典NP定义的操作不自足和概念错误（不可公度性丢失、对偶性消除、拷贝操作不可实现）。这与哥德尔不完备定理对希尔伯特第二问题的消解类似：算术一致性为真（在标准模型下），但无法在算术系统内部被证明。同样，$\mathbf{P} \neq \mathbf{NP}$为真，但无法在经典操作语义内部被证明。
	
	这一双重定位——``解决''（确认了$\mathbf{P} \neq \mathbf{NP}$的真值）与``消解''（揭示了经典框架的表达局限和概念错误）并存——精确地反映了本文工作的数学内容和哲学意义。它既不是一个空洞的``经典问题被解决''的声称，也不是一个悲观的``经典问题不可解''的断言，而是指出了一条中间道路：通过扩展计算模型的操作语义并修正经典定义的操作不自足和概念错误，使得分离成为可证的代数事实。
	
	\subsection{与哥德尔不完备定理的结构类比}
	
	本文的诊断与哥德尔不完备定理之间存在深刻的结构类比，这一类比有助于精确理解我们工作的性质。
	
	哥德尔第二不完备定理的论证结构是：（1）算术系统$S$，如果是一致的，那么它无法在$S$内部证明自身的一致性$\mathrm{Con}(S)$；（2）这不是说$\mathrm{Con}(S)$本身是假的——在标准模型下，算术是一致的，$\mathrm{Con}(S)$为真；（3）如果我们要证明$\mathrm{Con}(S)$，必须在比$S$更强的系统（如$S+\mathrm{Con}(S)$或$S+$超限归纳）中进行。
	
	我们的论证结构与此类似：（1）经典图灵机模型$M$，其操作语义无法内在地表达不可公度性（元定理），且经典NP定义在操作上不自足、在语义上残缺（丢失了不可公度性和对偶性）；（2）这不是说$\mathbf{P} \neq \mathbf{NP}$本身是假的——在集合论意义下，$\mathbf{P} \neq \mathbf{NP}$为真；（3）如果我们要证明$\mathbf{P} \neq \mathbf{NP}$，必须在比$M$语义更丰富的模型（CBTM/IVM）中进行，该模型修正了经典定义的操作不自足和概念错误。
	
	在这两种情况下，困境的根源都不是技术性的，而是\textbf{语言性的}。哥德尔的不可证命题在系统内部为真，但其证明需要比原系统更强的元理论工具。同样，$\mathbf{P} \neq \mathbf{NP}$在经典框架内为真，但其证明需要比经典模型更丰富的语义框架——一个能够内在地表达不可公度性、显式保留对偶性、并从根本上消除拷贝操作依赖的计算模型。
	
	当然，这一类比有其边界。哥德尔的不完备定理是关于\textbf{可证性}的——一致性不可证。我们的元定理是关于\textbf{可表达性}的——不可公度性不可内在表达。两者涉及不同的逻辑范畴。哥德尔的证明使用了自指性构造（哥德尔编码），而我们的证明使用了符号置换和语义重赋值。但这些技术差异不应掩盖两者共享的深层结构：两者都揭示了一个形式系统在表达和证明自身性质方面的根本局限。
	
	\section{结论与展望}
	
	本文在CBTM/IVM框架中完成了$\mathbf{P} \neq \mathbf{NP}$的严格证明。主要贡献包括：
	
	\begin{enumerate}
		\item \textbf{虚部验证机（IVM）与本质维度$\kappa(L)$}：构建了IVM作为CBTM的语义展开器，通过动态生成元管理机制解决了静态框架的维度坍缩和对偶性坍缩问题。定义了本质维度$\kappa(L)$，并严格证明了它的良定性、实现无关性和多项式上界。P类零维定理确立了所有$\mathbf{P}$语言的本质维度为零。分支-激活引理确立了维度消耗是IVM语义记账的强制性结果，独立于底层转移函数的设计，无法被任何算法绕过。IVM同时显式保留了每次分支的对偶性，修正了经典NP定义中对偶性被量词消除的错误，也从根本上消除了经典定义中隐含的虚部到实部拷贝操作。
		\item \textbf{障碍无关性论证}：系统证明了$\kappa(L)$不受相对化、自然证明和代数化三大障碍的约束。投影约束机制确保$\kappa(L)$在任何谕示下保持稳定；定义域限制使其规避了自然证明的``大尺度''条件；阈值判断的完全离散性使分支触发依赖于符号的正负号（对偶性），代数化谕示对此天然``失明''，基于敏感性的证明因此是谕示无关证明。
		\item \textbf{NP完全问题的维度下界}：通过信息论下界论证，证明了子集和问题的任何正确验证器都必须经历至少$n$次非确定性分支，结合分支-激活引理和素数平方根线性无关定理，得出$\kappa(\mathrm{SubsetSum}) = \Omega(n)$。验证敏感性定理独立证明了生成元激活是验证正确性的必要条件。维度保持归约定理将$\Omega(n)$的下界传递至所有$\mathbf{NP}$完全问题。
		\item \textbf{$\mathbf{P} \neq \mathbf{NP}$的证明}：通过反证法，假设$\mathbf{P} = \mathbf{NP}$导致子集和问题的本质维度同时为零和$\Omega(n)$的矛盾，故$\mathbf{P} \neq \mathbf{NP}$。整个证明在ZFC公理体系内完成。
		\item \textbf{经典框架的定位}：系统阐述了该证明与经典P vs NP问题的关系——$\mathbf{P}\neq\mathbf{NP}$在集合论意义下为真且在ZFC内可证，但该证明无法在经典操作语义内部被复制，因为经典NP定义在操作上不自足、在语义上残缺，使得本质维度$\kappa(L)$在经典模型中不可定义。
	\end{enumerate}
	
	该证明完全建立在障碍无关的不变量$\kappa(L)$之上，天然免疫于相对化、自然证明和代数化三大障碍。$\kappa(L)$作为一个基于语义的、障碍无关的不变量，为复杂度分类提供了全新的数学工具。它度量的是验证过程中必须激活的``代数维度''数目，这些维度由生成元的线性无关性刻画，是一种\textbf{语义层面的资源}。正是这种语义本质，使得$\kappa(L)$具备抵御外部干扰的内在能力，并能够捕捉到$\mathbf{P}$与$\mathbf{NP}$之间的本质差异。
	
	从更宏观的视角看，维度退化理论标志着计算复杂性研究的一次范式转换。旧范式（经典图灵机）的语言只描述符号的语法操作，无法触及计算背后的代数语义——不可公度性被外包，对偶性被消除，操作不自足被掩盖。新范式（CBTM/IVM）将非确定性的本质——不可公度性及其对偶结构——内建于操作语义，从根本上消除了拷贝操作的必要性，使P与NP的差异从模糊的哲学问题转化为精确的代数对比。这一转换可能具有超出P vs.\ NP问题的深远意义。其他复杂性类的分离问题——如$\mathbf{NP}$与$\mathrm{coNP}$、$\mathbf{P}$与$\mathrm{PSPACE}$——是否也存在类似的维度理论？随机性与非确定性之间的关系是否可以在代数维度谱系中获得新的解释？量子计算中的加速现象是否可以由量子本质维度$\kappa_q$来刻画？这些方向留待未来研究。
	
	\appendix
	
	\section{实现无关性的详细证明}
	
	\begin{proof}[定理\ref{thm:implementation-indep}的详细证明]
		设$M_1,M_2\in\mathcal{V}_L$。由IVM与NTM的等价性（见论文\cite{models2026}的参数化等价定理），存在多项式时间NTM $N_1,N_2$使得$L(N_i)=L$且$N_i$模拟$M_i$，反之亦然。由于$N_1$与$N_2$均为多项式时间NTM，标准模拟技术允许$N_1$在多项式时间内模拟$N_2$，且模拟过程中$N_2$的每一步至多被$N_1$的$q(|x|)$步模拟（$q$为某多项式）。提升到IVM层面，$M_1$模拟$M_2$的一步最多需要$q(|x|)$步，这期间可能触发分支并引入额外生成元，但引入数量不超过步数。因此对任何输入$x$，有$\kappa_{M_1}(x)\le \kappa_{M_2}(x)+q(|x|)$。对称地，反向不等式也成立（可能换一个多项式）。取$p=\max(q_1,q_2)$即得所需界限。
	\end{proof}
	
	\section{子集和维度记录的详细推导}
	
	\begin{proof}[定理\ref{thm:subset-sum-lower}中维度记录的详细推导]
		对任意子集和实例$(S,t)$，其中$S=\{s_1,\dots,s_n\}$，采用编码$\tau_i = s_i + 1\cdot\sqrt{p_i}$。IVM逐符号处理：
		\begin{itemize}
			\item 处理$\tau_1$：虚系数$1>\epsilon$，触发分支。根据分支-激活引理，IVM自动分配地址$\sqrt{2}$，同时对偶地引入$-\sqrt{2}$，并在激活路径上记录值$1$，生成元$\sqrt{2}$被激活。
			\item 依此类推，处理$\tau_i$时引入并激活$\sqrt{p_i}$，同时对偶地引入$-\sqrt{p_i}$。
			\item 由素数平方根线性无关定理（引理\ref{lem:commensurability}），这$n$个被激活的生成元线性无关，故$\kappa_M(x_n) = n$。
		\end{itemize}
		维度记录本身是分支-激活引理的直接结果，不依赖于元素$s_i$的具体取值。信息论下界论证（第5.1节）进一步确立了该下界对所有正确验证器的普适性：任何正确验证器都必须区分$2^n$种可能性，因此必须经历至少$n$次分支，每次分支必然激活一个生成元，故维度至少为$n$。
	\end{proof}
	
	\section{概念直觉：地址与值的比喻}
	
	本节提供一个直观的比喻，帮助读者理解CBTM到IVM语义展开的本质。此比喻仅为辅助理解之用，所有严格的数学定义和证明均已在正文中给出。
	
	\textbf{CBTM：地址即身份。} CBTM的地址空间固定为$\{\alpha, \beta\}$，大小恒为2。更准确地说，$\alpha$和$\beta$不仅是两个``地址''，更是两个\textbf{全局函数指针}，且互为对偶。每次非确定性分支的两条路径，分别调用这两个函数指针中的一个，调用$\alpha$意味着选择第一个选项并执行其关联的处理逻辑，调用$\beta$意味着选择第二个选项。无论发生多少次分支，机器始终只能在这两个共享的全局函数指针之间选择。后面的调用会覆盖之前调用的痕迹，因为纸带上同一个位置的符号只能携带一个虚部标记，新的写入会擦除旧的标记。因此，$m$次分支后，维度最多为$2$，对偶性也限于这一对指针。
	
	\textbf{IVM：地址上的值即身份。} IVM的地址空间是动态扩展的。每次非确定性分支触发时，系统会``开一个新页''，即分配一个全新的、与之前所有地址都线性无关的地址$\sqrt{p_i}$，同时对偶地引入$-\sqrt{p_i}$。两条路径的区别不再由``调用了哪个全局函数指针''决定，而是由该地址上写入的\textbf{值}决定：写入$1$（激活）意味着选择第一条路径（$+\sqrt{p_i}$），计算进入扩域，新的生成元$\sqrt{p_i}$被纳入代数结构；写入$0$（不激活）意味着选择第二条路径（$-\sqrt{p_i}$），计算停留在当前基域内，不引入新的生成元，但该地址空间仍被分配。$m$次分支后，地址空间扩展为$\{\sqrt{p_1}, \sqrt{p_2}, \dots, \sqrt{p_m}\}$，每条路径由一个$m$维布尔向量$\mathbf{a} \in \{0,1\}^m$唯一标记。维度等于被激活的地址数量，最多为$m$。对偶性也得以按分支步骤独立追踪——每次分支都有独立的对偶对$\sqrt{p_i}$与$-\sqrt{p_i}$。
	
	\textbf{核心转换：将地址转换为地址里的值。} CBTM到IVM的语义展开，其核心操作可以直观地理解为：
	\[
	\text{CBTM: 路径的标记} = \text{调用了哪个全局函数指针 } (\alpha \text{ 或 } \beta)
	\]
	\[
	\quad  \longrightarrow \quad \text{IVM: 路径的标记} = \text{在专属地址上写了什么值 } (1 \text{ 或 } 0).
	\]
	这一转换的数学基础是正文中定义的双射$\Phi: \{\alpha, \beta\}^m \to \{0,1\}^m$。
	
	\textbf{为什么这一转换是本质性的。}
	\begin{enumerate}
		\item \textbf{身份的解绑}：CBTM只有两个全局函数指针，被所有分支共享和复用，后面的调用会覆盖前面的痕迹。IVM为每次分支创建专属地址，每个地址上的标记一旦写入就不会被后续分支覆盖。
		\item \textbf{信息容量的解放}：CBTM只有两个函数指针，信息容量恒为$1$比特。IVM每次分支创建新地址，$m$次分支后信息容量为$m$比特，可以编码$2^m$种不同的路径标记。
		\item \textbf{维度的解放}：CBTM中维度受限于全局函数指针的数量（恒为$2$）。IVM中维度不再受限于此，而是由值为$1$的地址个数决定，最多可达$m$维。
		\item \textbf{对偶性的解放}：CBTM中所有分支共享同一对对偶生成元，对偶性坍缩为单一对偶对。IVM中每次分支有独立的对偶对$\sqrt{p_i}$与$-\sqrt{p_i}$，对偶性得以按分支步骤独立追踪。
		\item \textbf{从``标记''到``度量''}：CBTM只能标记，无法度量。IVM通过计算有多少个地址上的值为$1$，可以精确度量非确定性资源的消耗，每次激活都意味着一次真正的域扩张，消耗一个独立的代数维度。
		\item \textbf{不可归约性的保证}：在CBTM中，无法从代数上严格保证调用不同函数指针的两条路径不可相互归约。在IVM中，激活路径上的生成元与基域线性无关，从代数结构上确保了该路径无法被停留在基域内的确定性计算所模拟，这正是非确定性本质的形式化。
	\end{enumerate}
	
	简言之，\textbf{CBTM让所有非确定性选择共享同一对全局函数指针，导致维度坍缩和对偶性坍缩；IVM为每一次非确定性选择分配专属地址，以``是否扩域''的布尔值记录选择，使得维度能够线性累积，对偶性得以独立追踪。}这正是维度退化理论能够分离$\mathbf{P}$与$\mathbf{NP}$完全问题的核心机制。
	
	\bibliographystyle{plain}
	\bibliography{../cb}
	
\end{document}